\loeskopf{Einheit 4 von 5 · Nullstellen und Schnittpunkte der Parabel (Nr. 29--37)}

\erg{29}{a) $x^2 - 5x + 4 = 0$ \quad b) $x^2 + 2x = 3x + 6$ \quad
c) $x^2 - 4x = 12$ \quad d) $(x - 7)^2 - 9 = 0$ \quad e) $(x + 1)^2 = -x + 5$}

\erg{30}{a) $N_1(7 \mid 0)$, $N_2(-1 \mid 0)$ \quad b) $N_1(3 \mid 0)$,
$N_2(-7 \mid 0)$ \quad c) $N_1(7 \mid 0)$, $N_2(3 \mid 0)$ \quad
d) $N_1(-5 \mid 0)$, $N_2(-7 \mid 0)$ \quad e) $N(8 \mid 0)$, genau eine \quad
f) keine Nullstelle \quad g) $N_1(4 \mid 0)$, $N_2(-2 \mid 0)$ \quad
h) keine: Scheitel $S(-4 \mid 2)$ liegt über der $x$-Achse und die Parabel ist nach
oben geöffnet \quad i) zwei: Scheitel $S(3 \mid 7)$ liegt über der $x$-Achse und die
Parabel ist nach unten geöffnet}

\erg{31}{a) $x_1 = 5$, $x_2 = 2$ \quad b) $x_1 = -4$, $x_2 = -5$ \quad
c) $x_1 = 6$, $x_2 = -3$ \quad d) $x_1 = 4$, $x_2 = -6$ \quad
e) durch $2$: $x^2 - 5x + 6 = 0$; $x_1 = 3$, $x_2 = 2$ \quad
f) durch $3$: $x^2 + x - 6 = 0$; $x_1 = 2$, $x_2 = -3$ \quad
g) $x_{1,2} = 2 \pm \sqrt{3}$, also $3{,}73$ und $0{,}27$ \quad
h) $x_{1,2} = 4 \pm \sqrt{2}$, also $5{,}41$ und $2{,}59$}

\erg{32}{a) $N_1(-1 \mid 0)$, $N_2(3 \mid 0)$ \quad b) $S_1(1 \mid -4)$,
$S_2(4 \mid 5)$ \quad c) Die Aussage ist falsch. Eine waagerechte Gerade unterhalb des
Scheitels trifft die Parabel gar nicht, eine Gerade durch den Scheitel kann sie in
genau einem Punkt berühren.}

\erg{33}{a) $x^2 + 4x - 12 = 0$; $x_1 = 2$, $x_2 = -6$ \quad
b) $x^2 - 6x - 7 = 0$; $x_1 = 7$, $x_2 = -1$ \quad
c) $-x^2 - 2x + 15 = 0$, mal $(-1)$: $x^2 + 2x - 15 = 0$; $x_1 = 3$, $x_2 = -5$}

\erg{34}{a) $x^2 + 3x - 4 = 0$; $S_1(1 \mid 6)$, $S_2(-4 \mid -4)$ \quad
b) $x^2 - 3x - 10 = 0$; $S_1(5 \mid 21)$, $S_2(-2 \mid 0)$ \quad
c) $x^2 - 4x = 0$; $S_1(0 \mid 3)$, $S_2(4 \mid 27)$ \quad
d) $x^2 - 4x + 6 = 2x + 1$, also $x^2 - 6x + 5 = 0$; $S_1(1 \mid 3)$, $S_2(5 \mid 11)$
\quad e) $x^2 - 8x + 11 = -3x + 11$, also $x^2 - 5x = 0$; $S_1(0 \mid 11)$,
$S_2(5 \mid -4)$ \quad f) $x^2 - 6x + 8 = 0$; $S_1(2 \mid 3)$, $S_2(4 \mid 17)$}

\erg{35}{a) $p(3) = 9 - 5 = 4$ und $f(3) = 6 - 2 = 4$ -- der Punkt liegt auf beiden
Graphen. \quad b) ja: $p(-2) = 4 - 3 = 1$ und $f(-2) = 2 - 1 = 1$ \quad
c) nein: $p(1) = 1 + 3 - 4 = 0$, aber $f(1) = 2 - 1 = 1 \neq 0$; $R$ liegt nur auf der
Parabel}

\erg{36}{Fehler: Mia hat nur die $x$-Werte bestimmt und die $y$-Werte weggelassen --
$y = 0$ gilt für Nullstellen, nicht für Schnittpunkte mit einer Geraden. Richtig:
$f(1) = 5$ und $f(-4) = 0$, also $S_1(1 \mid 5)$ und $S_2(-4 \mid 0)$. \quad
a) $x^2 - 3x - 4 = 0$; $x_1 = 4$, $x_2 = -1$; $S_1(4 \mid 8)$, $S_2(-1 \mid 3)$}

\erg{37}{a) Auf der $x$-Achse ist jeder $y$-Wert null. Eine Nullstelle ist eine
Stelle mit $p(x) = 0$ -- genau dort trifft der Graph die $x$-Achse. \quad
b) Zum Beispiel $y = -5$: Die Parabel ist nach oben geöffnet, ihr kleinster Wert ist
$-4$. Alles unterhalb von $-4$ wird nie erreicht. \quad
c) Der Scheitel ist der höchste Punkt und liegt über der $x$-Achse. Von dort fällt der
Graph nach beiden Seiten beliebig weit -- also kreuzt er die $x$-Achse links und
rechts je einmal.}
